\clearpage
\begin{table}[H]\centering
\caption{Sampling variability with the fixed Hybrid checkpoint}\label{tab:S10}
\small
\begin{tabular}{p{.47\textwidth}rr}
\toprule
Summary & Internal & External \\
\midrule
WSIs & 49 & 168 \\
Additional bags per WSI & 20 & 20 \\
Median within-slide score SD & 0.0120 & 0.0150 \\
Within-slide score SD, interquartile range & 0.0019--0.0475 & 0.0065--0.0318 \\
Maximum within-slide score SD & 0.1297 & 0.1241 \\
Median within-slide score range & 0.0474 & 0.0581 \\
WSIs crossing 0.5 across the 20 bags & 7/49 & 18/168 \\
SS WSIs crossing 0.5 & 4/26 & 3/16 \\
MALT WSIs crossing 0.5 & 3/23 & 15/152 \\
Per-draw AUROC, mean (SD) & 0.9599 (0.0056) & 0.8749 (0.0097) \\
Per-draw AUROC, range & 0.9498--0.9682 & 0.8549--0.8943 \\
20-bag mean-score AUROC (secondary) & 0.9632 & 0.8775 \\
20-bag mean-score sensitivity (secondary) & 95.65\% & 92.76\% \\
20-bag mean-score specificity (secondary) & 80.77\% & 75.00\% \\
\bottomrule
\end{tabular}
\par\medskip\raggedright\footnotesize Each technical repeat uses 512 sampled feature rows and the same frozen checkpoint. Standard deviations use the sample definition. Crossing means that at least one score is below 0.5 and at least one is at or above 0.5 among the 20 additional bags. AUROC spread across draws is not a confidence interval. Mean-score results are secondary summaries across all 20 bags and do not replace the archived single-bag results. Source data: sampling\_summary.csv, sampling\_per\_slide.csv and sampling\_per\_draw\_metrics.csv.
\end{table}
\clearpage
\begin{table}[H]\centering
\caption{External model comparisons on shared reference bags}\label{tab:S11}
\small
\begin{tabular}{p{.18\textwidth}p{.23\textwidth}p{.23\textwidth}p{.23\textwidth}}
\toprule
Model & AUROC (95\% interval) & Sensitivity (95\% interval) & Specificity (95\% interval) \\
\midrule
Hybrid & \textbf{0.8721} (0.7294--0.9737) & \textbf{93.42\%} (88.31--96.39) & \underline{68.75\%} (44.40--85.84) \\
ABMIL & \underline{0.7903} (0.6275--0.9268) & 42.76\% (35.17--50.71) & \textbf{81.25\%} (56.99--93.41) \\
CLAM-SB-style & 0.6591 (0.4568--0.8491) & 58.55\% (50.60--66.08) & 62.50\% (38.64--81.52) \\
RRTMIL-style & 0.7743 (0.6114--0.9100) & \underline{89.47\%} (83.59--93.42) & 50.00\% (28.00--72.00) \\
\bottomrule
\end{tabular}
\medskip
\begin{tabular}{p{.32\textwidth}rp{.27\textwidth}}
\toprule
Paired difference & AUROC difference & 95\% interval \\
\midrule
Hybrid minus ABMIL & 0.0818 & -0.0136 to 0.2011 \\
Hybrid minus CLAM-SB-style & 0.2130 & 0.0843 to 0.3528 \\
Hybrid minus RRTMIL-style & 0.0979 & 0.0415 to 0.1649 \\
\bottomrule
\end{tabular}
\par\medskip\raggedright\footnotesize Bold/underline mark the highest/second-highest point estimate in each metric column. All four original checkpoints were evaluated on identical per-slide reference bags from 168 WSIs (152 MALT and 16 SS). Sensitivity and specificity use threshold 0.5 and Wilson intervals. AUROC and paired-difference intervals in this new analysis use 10,000 label-stratified WSI bootstrap draws; the archived Hybrid intervals in main Table 3 use 2,000 draws. Differences are unadjusted descriptive comparisons of complete model implementations, not isolated aggregation mechanisms or confirmatory family-wise inference. Source data: reference\_model\_metrics.csv and external\_paired\_comparisons.csv.
\end{table}
\clearpage
\begin{table}[H]\centering
\caption{Archived segmentation-contour coverage of the eight displayed regions}\label{tab:S12}
\small Individual-case rows are omitted from this public repository copy.
\end{table}
\clearpage
\begin{table}[H]\centering
\caption{Descriptive Brier scores of archived predictions}\label{tab:S13}
\small
\begin{tabular}{llrl}
\toprule
Cohort & Model & WSIs & Brier score (95\% interval) \\
\midrule
Internal & Hybrid & 49 & 0.0895 (0.0429--0.1463) \\
Internal & ABMIL & 49 & 0.1056 (0.0396--0.1827) \\
Internal & CLAM-SB-style & 49 & 0.1002 (0.0561--0.1525) \\
Internal & RRTMIL-style & 49 & \underline{0.0894} (0.0341--0.1569) \\
Internal & Three-model ensemble & 49 & \textbf{0.0829} (0.0412--0.1330) \\
External & Hybrid & 168 & 0.0919 (0.0702--0.1156) \\
\bottomrule
\end{tabular}
\par\medskip\raggedright\footnotesize Bold/underline mark the lowest/second-lowest point estimate within each cohort; a lone estimate is unranked. Brier score is the mean squared difference between the uncalibrated sigmoid score and the binary label (MALT = 1). Intervals use 10,000 label-stratified WSI bootstrap resamples of fixed predictions. Seeds are recorded in the source table. No calibration function was fitted to either evaluation cohort. Brier score combines several aspects of probabilistic prediction and does not independently isolate calibration or establish clinical utility. Source data: brier\_descriptive.csv.
\end{table}
\clearpage
\begin{table}[H]\centering
\caption{Dimensions and parameter counts of the evaluated complete models}\label{tab:S14}
\small
\begin{tabular}{lrrrrr}
\toprule
Model & Head input & Attention width & Parameters & Epoch budget & Batch \\
\midrule
Hybrid & 256 & 256 & 1,521,672 & 30 & 8 \\
ABMIL & 1536 & 512 & 6,690,306 & 50 & 16 \\
CLAM-SB-style & 256 & 256 & 591,874 & 20 & 8 \\
RRTMIL-style & 256 & 256 & 1,711,874 & 20 & 4 \\
\bottomrule
\end{tabular}
\par\medskip\raggedright\footnotesize All models receive 1,536-dimensional UNI2-h features and use 512 sampled rows. Head input is the dimension passed from the slide encoder to the common classifier construction; its hidden width is 256. ABMIL retains 1,536-dimensional patch features with additional transformations before and after attention pooling, whereas Hybrid projects features to 256 dimensions and fuses five vectors with eight statistics. These differences, along with training budgets, prevent attribution to pooling alone. Parameter counts were obtained by strict loading of the archived checkpoints; they do not measure runtime, memory, or energy. Source data: model\_checkpoint\_inventory.csv and checkpoint\_tensor\_shapes.json.
\end{table}
